\subsection{Setup}
\label{sec:exp-setup}
\paragraph{Data and ground truth.}
We explain action recognition on MSR-Action3D~\cite{li2010msraction3d}: 567 depth videos of 20 actions by 10 subjects, with the standard split (subjects 1--5 for training, 270 videos; 6--10 for testing, 297 videos). Depth maps become point-cloud videos with the MeteorNet protocol~\cite{liu2019meteornet} used by all three official backbones; we use centred 24-frame clips with 2048 points per frame, re-sampled independently in every frame (no point identity). MSR-Action3D also ships 20-joint skeletons. We lift them into the same coordinates, where a median \NumGTFitMedian\% of the points lie within about 12\,cm of a bone, and label every point with its nearest bone. Some videos store no skeleton for a run of frames; every ground-truth-based number leaves out the \NumGTMissing\ benchmark clips concerned. Points of bones that travel at least half as far as the clip's most active bone form the \emph{moving} ground truth (on average \NumGTMovingPct\% of a clip's points); points of bones that travel at most a fifth as far form the \emph{static} ground truth (\NumGTStaticPct\%; Appendix~\ref{supp:impl}). The main benchmark explains every test video that is at least 24 frames long (\NumPFTClips\ clips; \NumGTValid\ with a complete skeleton). The cross-backbone study, the ablations and the stability study use the first two test clips of every action (\NumCrossClips\ clips).

\paragraph{Backbones.}
We explain the \emph{unmodified official code} of the convolutional PSTNet~\cite{fan2021pstnet}, the attention-based P4Transformer~\cite{fan2021p4transformer} and the attention-free state-space Mamba4D~\cite{liu2025mamba4d}. No pretrained weights exist, so we train each from scratch at its official width on one A100 GPU, replacing the CUDA-only operators by numerically equivalent PyTorch versions (Appendix~\ref{supp:impl}). Test accuracies (video level) are \BBaccPFT/\BBaccPST/\BBaccMBA\,\% for P4Transformer/PSTNet/Mamba4D (published: 90.9/91.2/92.7). A second, laptop-scale training (reduced width, 1024 points) supports the same conclusions; Appendix~\ref{supp:scale} compares the two.

\paragraph{Baselines.}
\emph{Floors:} random; geometry-edge (local point spacing); motion prior (model-blind nearest-neighbour flow between frames). \emph{Gradient:} Grad$\times$Input, SmoothGrad~\cite{smilkov2017smoothgrad}, Integrated Gradients~\cite{sundararajan2017ig}, and 3D Grad-CAM~\cite{selvaraju2017gradcam} on the tube-embedding tokens. \emph{Perturbation:} region occlusion (deleting each tracked region, as in point-cloud saliency~\cite{zheng2019pointcloudsaliency}), OccAM~\cite{schinagl2022occam}, and extremal perturbations~\cite{fong2019extremal}. \emph{Attention (P4Transformer only):} rollout~\cite{abnar2020rollout}, transformer attribution~\cite{chefer2021transformer}, and raw attention: the attention each token receives in the last layer, used \emph{as} the explanation.
\paragraph{\method.} $m{=}12$ tracked regions; local estimator (37 passes) and context estimator ($S{=}256$, $q{=}\tfrac12$). An attribution map gives each point its region's $\Nec$; localisation uses $\Smot$ (moving GT) and $\Sgeo$ (static GT).

\paragraph{Metrics.}
\emph{Faithfulness} (no GT): deletion/insertion AUC of the target-class probability~\cite{petsiuk2018rise}, comprehensiveness/sufficiency (AOPC over $\{1,5,10,20,50\}\%$)~\cite{deyoung2020eraser}, and Sensitivity-$n$~\cite{ancona2018sensitivity}; a deleted point is removed and its frame re-sampled. \emph{Localisation} (skeleton GT): Pointing Game~\cite{zhang2018pointing} and energy-based EBPG~\cite{wang2020scorecam}. \emph{Sanity:} cascading model randomisation~\cite{adebayo2018sanity}. \emph{Statistics.} Every score is computed per clip and reported as the mean over clips. A $p$-value compares two methods on the same clips with a two-sided paired Wilcoxon signed-rank test~\cite{wilcoxon1945}; $p$-values in the text are uncorrected, and a $^\ast$ in a table marks \method\ better than the strongest baseline after Bonferroni correction over that table's metrics~\cite{dunn1961}. Intervals are 95\% percentile bootstrap intervals over clips (2000 resamples)~\cite{efron1979bootstrap}. Appendix~\ref{supp:provenance} lists the data, models and code behind every table and figure.
